\chapter{Experiments}

\section{Reading Rule}

Each experiment should be read as computational evidence with an explicit
boundary. A result file may report a successful run, but the corresponding
mathematical claim still needs independent review.

\section{Experiment 001: Square-Zero Sanity Checks}

Experiment 001 has status \emph{run}. It checks one-dimensional split
square-zero extensions \(E = J \ltimes M\) with exact rational arithmetic.

The result is a sanity check for scalar actions in the idempotent and
zero-multiplication one-dimensional cases. It does not prove the Beck-module
classification and does not change the status of \path{CLAIM-0003}.

\section{Experiment 002: Two-Dimensional Jordan Algebras}

Experiment 002 has status \emph{not run yet}. Its goal is to enumerate or
import small two-dimensional Jordan algebra examples and test derivations,
indecomposables, and candidate differential modules.

\section{Experiment 003: Beck Module Identities}

Experiment 003 has status \emph{not run yet}. It is intended to symbolically
derive the identities imposed on \(M\) by the Jordan identity on a square-zero
extension \(J \ltimes M\).

\section{Experiment 004: TKK Comparison}

Experiment 004 has status \emph{not run yet}. Its purpose is to compare
intrinsic Jordan Quillen invariants with Lie homology of the associated TKK
Lie algebra on small examples.

\section{Experiment 005: Derived Indecomposables Toy Benchmark}

Experiment 005 has status \emph{run}. It benchmarks a one-generator toy model
for derived indecomposables on \(F(x)\), \(F(x)/(x^2-x)\), \(F(x)/(x^2)\),
\(F(x)/(x^3)\), and \(F(x)/(x^4)\).

The rule is that after applying \(Q\), only the linear part of a relation
survives. The current generated result says that all five planned examples
match their expected \(H_0\) and \(H_1\) dimensions.

\section{\texorpdfstring{Experiment 006: Presentation-Invariance Stress Test For \(J_3\)}{Experiment 006: Presentation-Invariance Stress Test For J3}}

Experiment 006 has status \emph{run}. It studies
\[
  J_3 = (t)/(t^3)
\]
through Presentation A, \(F(x)/(x^3)\), and a naive diagnostic for Presentation
B, \(F(x,y)/(y-x^2, xy, y^2)\).

The current result computes \(H_0=1, H_1=1\) for Presentation A. The naive
one-step diagnostic for Presentation B computes \(H_0=1, H_1=2\), exposing an
extra false class caused by the incomplete relation complex.

\section{\texorpdfstring{Experiment 007: Two-Step Cofibrant Replacement Toy For \(J_3\)}{Experiment 007: Two-Step Cofibrant Replacement Toy For J3}}

Experiment 007 has status \emph{run}. It keeps Presentation B from Experiment
006,
\[
  F(x,y)/(y-x^2, xy, y^2),
\]
and adds one explicit degree \(2\) generator \(b\) to kill the extra naive
class represented by \(a_3\).

The hand-coded two-step replacement uses degree \(0\) generators \(x,y\),
degree \(1\) generators \(a_1,a_2,a_3\), and differential
\[
\begin{aligned}
  d(a_1) &= y-x^2,\\
  d(a_2) &= xy,\\
  d(a_3) &= y^2,\\
  d(b) &= a_3 - x a_2 - (y+x^2)a_1 + x(xa_1).
\end{aligned}
\]
After applying indecomposables \(Q(-)=(-)/(-)^2\), the complex is
\[
  Q_2=\langle b\rangle \longrightarrow
  Q_1=\langle a_1,a_2,a_3\rangle \longrightarrow
  Q_0=\langle x,y\rangle
\]
with
\[
  d_Q(a_1)=y,\qquad d_Q(a_2)=0,\qquad d_Q(a_3)=0,\qquad d_Q(b)=a_3.
\]

The generated result records \(d_Q^2=0\), and the symbolic check behind the
degree \(2\) attachment reduces in this fixed one-generated example to
\[
  x^2x^2 - x(xx^2)=0,
\]
using power-associativity in the subalgebra generated by \(x\).

The computed homology dimensions are
\[
  H_0=1,\qquad H_1=1,\qquad H_2=0,
\]
with representatives \(x\) in \(H_0\), \(a_2\) in \(H_1\), and no \(H_2\)
representative. Compared with Experiment 006, the naive Presentation B
\(H_1\)-dimension drops from \(2\) to \(1\).

The boundary is essential: this verifies the fixed low-degree two-step repair
for \(J_3\) Presentation B through degree \(1\). It is not an automatic syzygy
search, not a general cofibrant replacement algorithm, and not a proof of full
presentation invariance.
